% ============================================================
% OP-Q11 CONVERSE CALCULATION — momentum-space OS test
% Version: v3 (2026-06-10). STATUS: cleanup round per GPT review
% of v2. Changes vs v2: (1) Family II index bookkeeping rewritten
% with a single consistent x,y convention and explicit reduction
% to the standard fixed-momentum form; (2) convention-check
% paragraph added (Sec. 0.2); (3) Majorana/charge-conjugation
% translation caveat recorded. No new physics relative to v2.
% Confirmed by review (conditionally, in trial convention):
%   M_I=(omega+h_+)/2omega   projector  (Family I, Grassmann)
%   M_II^Gr=(omega-h_-)/2omega projector (Family II, Grassmann)
%   M_II^com=(h_- -omega)/2omega, spectrum {0,-1} (commuting)
% ============================================================

\section*{Momentum-space OS test for the free minimal Euclidean
Dirac kernel (v3)}

% ------------------------------------------------------------
\subsection*{0. Conventions (trial set, complete)}
Euclidean $x=(w,\vec x)$, reflection
$\Theta:(w,\vec x)\to(-w,\vec x)$, consistent with $t=w$.
Hermitian gammas $\{\gamma_A,\gamma_B\}=2\delta_{AB}$.
Mass gap $m>0$ throughout.
Minimal kernel $D_E=\gamma_A\partial_A+m$,
$\tilde S_E(p)=\dfrac{m-i\gamma_Ap_A}{p^2+m^2}$,
Berezin pairing
$\langle\psi_\alpha(x)\bar\psi_\beta(y)\rangle
=S_{E,\alpha\beta}(x-y)$.
Reordering rule (the only entry point of statistics):
$\langle\bar\psi_\beta(y)\psi_\alpha(x)\rangle
=\sigma S_{E,\alpha\beta}(x-y)$,
$\sigma=-1$ (Grassmann), $\sigma=+1$ (commuting).
Reflection action (trial):
$\Theta\psi(x)=\bar\psi(\Theta x)\gamma_0$,
$\Theta\bar\psi(x)=\gamma_0\psi(\Theta x)$,
antilinear, order-reversing.
OS form: $(F,G)\equiv\langle F\cdot\Theta G\rangle$.
Fourier convention:
$g(w,\vec x)=\int\!\frac{d^3p}{(2\pi)^3}\,
e^{i\vec p\cdot\vec x}\,\tilde g(w,\vec p)$.
Notation: $\omega=\sqrt{\vec p^{\,2}+m^2}$,
$h_\pm(\vec p)=m\gamma_0\pm i\gamma_0\vec\gamma\cdot\vec p$,
hermitian, $h_\pm^2=\omega^2$, so
$(\omega\pm h_\pm)/2\omega$ are orthogonal rank-2 projectors.

\subsubsection*{0.1 Order convention}
At the quadratic level the alternative assignment
$(F,G)=\langle\Theta F\cdot G\rangle$ differs by the global
reordering sign $\sigma$; for Grassmann fields exactly one
assignment is OS-positive.
The order above is fixed by the internal consistency requirement
that \emph{both} elementary families below come out positive in the
Grassmann case; this removes the v1 sign ambiguity.

\subsubsection*{0.2 Convention check (recorded; final
reconciliation pending)}
The trial set above is constructed to lie in the standard Euclidean
Fermi-field OS convention family: reflection through the Euclidean
time gamma, antilinear and order-reversing on the field algebra,
with the sesquilinear form positive on the half-space algebra
(Osterwalder--Schrader Euclidean Fermi fields; L\"uscher's
transfer-matrix construction; Montvay--M\"unster \S4, with
$\gamma_4\leftrightarrow\gamma_0$ renaming).
A line-by-line reconciliation against these sources, and against
the Euclidean Dirac block ultimately adopted in the main text
(the current Lorentzian Dirac section, \S\ref{sec:dirac_eq}, fixes no Euclidean conventions), is required at
implementation time and is the remaining condition on Q11.1a.
If the main ECT spinor sector is later written in
Majorana/charge-conjugated variables, the calculation must be
translated into that convention; the natural $\psi\bar\psi$
ordering used here makes the charge-conjugation matrix unnecessary
for the present proof route.

% ------------------------------------------------------------
\subsection*{1. Contour lemma}
For $W>0$, $s=\pm$:
\[
\int\frac{dp_0}{2\pi}\,
\frac{e^{\,is\,p_0W}\,(m-i\gamma_0p_0-i\vec\gamma\cdot\vec p)}
{p_0^2+\omega^2}
=\frac{e^{-\omega W}}{2\omega}\,
\bigl(m+s\,\omega\gamma_0-i\vec\gamma\cdot\vec p\bigr).
\]

% ------------------------------------------------------------
\subsection*{2. Grassmann case ($\sigma=-1$): both families
reflection positive}
\emph{Family I.}
$F=\int_{w>0}d^4x\,f^\dagger(x)\psi(x)$,
$\Theta F=\int_{w>0}d^4y\,\bar\psi(\Theta y)\gamma_0 f(y)$,
\[
Q_{\rm Gr}^{\rm I}[f]
=\langle F\,\Theta F\rangle
=\int\!\!\int d^4x\,d^4y\;
f^\dagger(x)\,S_E(x-\Theta y)\,\gamma_0\,f(y),
\]
no reordering needed.
With $x-\Theta y=(w_x+w_y,\vec x-\vec y)$ and the contour lemma
($s=+$), the $w$-dependence factorises through
$h_f(\vec p)=\int_0^\infty dw\,e^{-\omega w}\tilde f(w,\vec p)$ and
\[
Q_{\rm Gr}^{\rm I}[f]
=\int\frac{d^3p}{(2\pi)^3}\;
h_f^\dagger(\vec p)\,\mathcal M_{\rm I}(\vec p)\,h_f(\vec p),
\qquad
\mathcal M_{\rm I}
=\frac{(m+\omega\gamma_0-i\vec\gamma\cdot\vec p)\gamma_0}{2\omega}
=\frac{\omega+h_+(\vec p)}{2\omega}\;\ge0 .
\]

\emph{Family II (rewritten; single $x,y$ convention).}
Define
\[
G=\int_{w>0}d^4x\;\bar\psi(x)\,g(x),
\qquad
\Theta G=\int_{w>0}d^4y\;g^\dagger(y)\,\gamma_0\,\psi(\Theta y).
\]
Then, in components,
\[
G\cdot\Theta G
=\int\!\!\int d^4x\,d^4y\;
g_\alpha(x)\,g^{*}_{\sigma'}(y)\,(\gamma_0)_{\sigma'\beta}\;
\bar\psi_\alpha(x)\,\psi_\beta(\Theta y),
\]
and taking the expectation with the reordering rule of \S0,
$\langle\bar\psi_\alpha(x)\psi_\beta(\Theta y)\rangle
=\sigma\,S_{E,\beta\alpha}(\Theta y-x)$:
\[
Q^{\rm II}[g]
=\sigma\int\!\!\int d^4x\,d^4y\;
g^\dagger(y)\;\gamma_0\,S_E(\Theta y-x)\;g(x).
\]
Here $\Theta y-x=(-(w_x+w_y),\vec y-\vec x)$, so the contour lemma
applies with $s=-$ and momentum $\vec q$ conjugate to
$\vec y-\vec x$:
\[
\gamma_0\,S_E(\Theta y-x)\Big|_{\vec q}
=e^{-\omega W}\,
\frac{m\gamma_0-\omega-i\gamma_0\vec\gamma\cdot\vec q}{2\omega}
=e^{-\omega W}\,\frac{h_-(\vec q)-\omega}{2\omega}.
\]
After the Fourier transform and relabelling of dummy variables
(the $e^{-\omega w_x}$, $e^{-\omega w_y}$ factors attach to $g(x)$
and $g^\dagger(y)$ respectively, producing the same half-space
profile $h_g$), the form reduces to the standard fixed-momentum
expression
\[
Q^{\rm II}[g]
=\sigma\int\frac{d^3q}{(2\pi)^3}\;
h_g^\dagger(\vec q)\;
\frac{h_-(\vec q)-\omega}{2\omega}\;
h_g(\vec q),
\qquad
h_g(\vec q)=\int_0^\infty dw\,e^{-\omega w}\,\tilde g(w,\vec q).
\]
For Grassmann fields ($\sigma=-1$):
\[
\mathcal M_{\rm II}^{\rm Gr}(\vec q)
=\frac{\omega-h_-(\vec q)}{2\omega}\;\ge0
\quad\text{(orthogonal projector)}
\;\Longrightarrow\;
Q_{\rm Gr}^{\rm II}[g]\ge0 .
\]
[The spatial-momentum orientation differs between the families
($h_+$ vs $h_-$); both conclusions are orientation-independent:
positivity is the projector property, and the converse below uses
only $\mathrm{spec}\,h_-=\{\pm\omega\}$.]
Extension from the two elementary quadratic families to the full
graded polynomial algebra is the standard OS fermionic construction
(cited, not re-derived).

% ------------------------------------------------------------
\subsection*{3. Commuting case ($\sigma=+1$): explicit negative
modes}
The strict negative-mode test assumes $m>0$; the massless case
requires separate infrared treatment and is not part of Q11.1a.
Family II with $\sigma=+1$ gives
\[
\mathcal M_{\rm II}^{\rm com}(\vec q)
=\frac{h_-(\vec q)-\omega}{2\omega},
\qquad
\mathrm{spec}=\Bigl\{\,0,\;-1\,\Bigr\}\ \text{(each twice)} .
\]
\emph{Explicit negative-norm test function.}
Take $\vec q\approx0$, $u_-$ with $\gamma_0u_-=-u_-$
(so $h_-(0)u_-=-\omega u_-$), and
$g(w,\vec x)=\chi(w)\rho(\vec x)u_-$ with $\chi$ supported in
$w>0$, $\rho$ a broad envelope:
\[
Q_{\rm com}[g]\;\longrightarrow\;-\,|h_g(0)|^2\;<\;0 .
\]
Family I remains positive for $\sigma=+1$ (no reordering occurs
there); reflection positivity fails because the OS form must be
non-negative on \emph{all} test functionals.
\emph{Mechanism.}
The statistics assignment enters in exactly one place: the
reordering sign $\sigma$.
Grassmann statistics converts the Family-II kernel into the
projector $(\omega-h_-)/2\omega$; commuting statistics leaves
$(h_--\omega)/2\omega$ with spectrum $\{0,-1\}$.

\emph{Diagnostic remark (not part of the proof).}
A separate consistency route would require constructing a positive
commuting Gaussian covariance with spinor indices; if enforcing
ordinary symmetric covariance removes the first-order Dirac
numerator, the resulting object is no longer the same minimal
kernel.
Diagnostic only; the converse rests on the negative-mode test
above.

% ------------------------------------------------------------
\subsection*{4. Status and remaining conditions}
\begin{enumerate}[label=(\alph*)]
  \item Review status: Families I/II (Grassmann) and the
    $\{0,-1\}$ commuting spectrum are confirmed \emph{conditionally
    in the trial convention} (review round v2).
  \item Remaining condition on Q11.1a: the final convention check
    of \S0.2 (line-by-line reconciliation with the OS/Fermi
    literature and with the Euclidean Dirac block to be adopted in
    the main text).
  \item Safe status wording: ``candidate Level~A mathematical
    result for a free minimal Euclidean Dirac kernel, with the
    momentum-space converse calculation completed in trial OS
    convention; pending final convention check.''
  \item Scope: minimal first-order kernel; $m>0$; pointlike
    primary-branch excitations; no higher-derivative kernels;
    Majorana/charge-conjugated reformulation would require
    translation (\S0.2).
  \item Forbidden wording remains in force: not ``ECT derives
    spin-statistics''; correct: ``the free-kernel OS
    spin--statistics test has a completed candidate calculation;
    its application to ECT fermions remains conditional on
    deriving or verifying the ECT spinor kernel and spinorial
    reflection-positive OS structure.''
\end{enumerate}
% ============================================================
